%%% Ví dụ: Hóa học lớp 9 - Chủ đề Oxi hóa khử %%%
%%% Gồm đầy đủ 4 loại câu hỏi %%%

%%% LOẠI 1: Trắc nghiệm nhiều lựa chọn %%%
%%%=============EX_1=============%%%
\begin{ex}%[9K8NB1-1]
	Phản ứng oxi hóa - khử là phản ứng hóa học trong đó có sự thay đổi đại lượng nào sau đây của các nguyên tử?
	\choice
	{Số khối}
	{\True Số oxi hóa}
	{Số electron lớp ngoài cùng}
	{Số nơtron trong hạt nhân}
	\loigiai{
		Phản ứng oxi hóa - khử là phản ứng hóa học trong đó có sự thay đổi số oxi hóa của một số nguyên tố. Quá trình oxi hóa là sự tăng số oxi hóa, còn quá trình khử là sự giảm số oxi hóa.
	}
\end{ex}

%%%=============EX_2=============%%%
\begin{ex}%[9K8NB2-1]
	Trong phản ứng: $Fe + CuSO_4 \xrightarrow FeSO_4 + Cu$, chất khử là
	\choice
	{$CuSO_4$}
	{$Cu$}
	{\True $Fe$}
	{$FeSO_4$}
	\loigiai{
		Trong phản ứng trên, $Fe$ có số oxi hóa tăng từ $0$ lên $+2$ (nhường electron), nên $Fe$ là chất khử.
	}
\end{ex}

%%% LOẠI 2: Trắc nghiệm đúng sai %%%
%%%=============TF_1=============%%%
\begin{ex}%[9K8TH1-2]
	Xét phản ứng tạo thành Calcium chloride từ đơn chất: $Ca + Cl_2 \xrightarrow CaCl_2$. Đánh giá tính đúng/sai của các phát biểu sau:
	\choiceTF
	{\True Trong phản ứng trên thì mỗi nguyên tử Calcium nhường 2e}
	{Số oxi hóa của Ca và Cl trước phản ứng lần lượt là +2 và -1}
	{Nếu dùng 4 gam Calcium thì số mol electron Chlorine nhận là 0,4 mol}
	{\True Liên kết trong phân tử $CaCl_2$ là liên kết ion}
	\loigiai{
		Phản ứng: $\overset{0}{Ca} + \overset{0}{Cl}_2 \xrightarrow \overset{+2}{Ca}\overset{-1}{Cl}_2$
		\begin{itemchoice}[T1,F2,F3,T4]
			\itemch Quá trình oxi hóa Ca: $Ca \xrightarrow Ca^{2+} + 2e^-$. Mỗi nguyên tử Ca nhường 2 electron. Phát biểu đúng
			\itemch Trước phản ứng, $Ca$ và $Cl_2$ là các đơn chất nên có số oxi hóa bằng $0$. Phát biểu sai
			\itemch $n_{Ca} = \dfrac{4}{40} = 0{,}1$ mol. Số mol electron Ca nhường: $n_e = 0{,}1 \times 2 = 0{,}2$ mol. Theo bảo toàn electron, $Cl_2$ nhận $0{,}2$ mol electron. Phát biểu sai
			\itemch $CaCl_2$ tạo từ kim loại điển hình và phi kim điển hình, hiệu độ âm điện lớn nên là liên kết ion. Phát biểu đúng
		\end{itemchoice}
	}
\end{ex}

%%% LOẠI 3: Trả lời ngắn %%%
%%%=============SA_1=============%%%
\begin{ex}%[9K8VD1-1]
	Cho phản ứng: $2Al + 3H_2SO_4 \xrightarrow Al_2(SO_4)_3 + 3H_2$. Tính tổng hệ số cân bằng (tổng các hệ số nguyên tối giản) của phản ứng trên.
	\shortans{$9$}
	\loigiai{
		Các hệ số cân bằng: $2 + 3 + 1 + 3 = 9$.
	}
\end{ex}

%%% LOẠI 4: Tự luận %%%
%%%=============BT_1=============%%%
\begin{bt}%[9K8VD2-1]
	Cho $5{,}6$ gam sắt ($Fe$) tác dụng với dung dịch đồng(II) sunfat ($CuSO_4$) dư.
	\begin{enumerate}
		\item Viết phương trình hóa học của phản ứng và xác định chất khử, chất oxi hóa, quá trình khử, quá trình oxi hóa.
		\item Tính khối lượng đồng ($Cu$) thu được sau phản ứng.
	\end{enumerate}
	\loigiai{
		\textbf{1.} Phương trình hóa học:
		\[ Fe + CuSO_4 \xrightarrow FeSO_4 + Cu \]
		\begin{itemize}
			\item Chất khử: $Fe$ (số oxi hóa tăng từ $0$ lên $+2$)
			\item Chất oxi hóa: $CuSO_4$ (Cu có số oxi hóa giảm từ $+2$ xuống $0$)
			\item Quá trình oxi hóa: $Fe \xrightarrow Fe^{2+} + 2e$
			\item Quá trình khử: $Cu^{2+} + 2e \xrightarrow Cu$
		\end{itemize}

		\textbf{2.} Tính khối lượng Cu thu được:
		\begin{eqnarray*}
			n_{Fe} &=& \dfrac{5{,}6}{56} = 0{,}1 \text{ mol} \\
			n_{Cu} &=& n_{Fe} = 0{,}1 \text{ mol} \\
			m_{Cu} &=& 0{,}1 \times 64 = 6{,}4 \text{ g}
		\end{eqnarray*}
		Vậy khối lượng đồng thu được là $6{,}4$ gam.
	}
\end{bt}
